\subsection{实验结论}
实验结果表明,即使只是对向量求和程序进行较为基础的代码优化,如减少过程调用、直接访问数组元素、使用局部变量累加等,也能够显著提升程序运行效率,充分体现了程序组织方式和访存方式对性能的直接影响。

在多种优化方法中,循环展开和多累加器设计是提升向量求和性能的核心手段。其中取得最佳优化效果的是优化四,即8路累加器与8路循环展开版本。该版本在保证计算结果正确的前提下,将核心求和过程的周期数由基线版本的\texttt{50012}降低到\texttt{22556},性能提升约为\textbf{54.9\%}。

同时,实验也表明并非所有“更底层”的写法都一定更优。例如,优化三虽然尝试将数组下标访问改为指针遍历,但实际性能反而略低于优化二。这说明程序优化必须结合具体平台、编译器和运行环境,通过实验数据进行验证,而不能仅凭理论推测判断。

此外,在优化四基础上增加函数级\texttt{O3}编译器优化属性后,程序性能未出现进一步明显提升。这说明当前优化四代码本身已经较适合编译器自动优化,其结构已较接近当前实验环境下的高效实现形式。因此,本实验中的主要性能收益仍然来自手工代码优化,而编译器优化在该阶段未表现出额外显著作用。

总体来看,香山仿真环境和NEMU模拟器能够有效支持任务二中向量求和程序的实现、优化与性能分析。虽然仿真环境下的绝对性能会受到软件模拟开销影响,但其对于验证优化思路、观察性能变化趋势以及分析不同实现方式之间的相对效果仍然具有重要意义。